\chapter{气体分子动理论}

\section{物质的性质}

从本章起我们将开始一个新的课题，这个课题将占用相当的时间。它是从物质由大量原子或基本单元组成，它们之间存在着电相互作用，并且从遵循力学定律这种物理观点出发，对物质的性质进行分析的第一部分。我们企图了解为什么不同的原子集合会表现出它们所具有的特色。

显然这是一个困难的课题，我们要在一开始就强调指出，事实上这个课题是非常困难的，而且所用的处理方法也必须与迄今为止我们处理其他课题的方法都不相同。在力学和光学中，我们能够从某些定律，比如牛顿定律，或者加速电荷所产生的电场的公式的严格叙述开始。由此可以从本质上理解一大批现象，并且从此以后，这些定律就为我们理解力学与光学提供了基础。这就是说我们在以后可以学到更多的东西，但是我们并不是学习不同的物理内容，而只是学习用以处理问题的更好的数学分析方法。

在研究物质的性质时，我们不能有效地沿用这条途径，只能以一种非常初步的方式来讨论物质；直接从具体的基本定律出发来分析这样的课题是非常复杂的，因为这些定律无非是力学与电学的定律。但它们与我们希望研究的物质的性质之间相隔太远；从牛顿定律得出物质的性质要经过非常多的步骤，而这些步骤本身也是相当复杂的。我们现在开始采取其中的某些步骤，虽然我们的许多分析相当精确，但所得出的最后结果却越来越不精确。对于物质的性质，我们只能有一个大致的了解。

我们之所以只能用如此不完善的方式来进行分析，其理由之一是它在数学上要求对概率论具有深刻的理解；我们并不要求知道每个原子实际上在哪里运动，而是要求知道平均说来有多少原子在这里运动，有多少原子在那里运动，产生各种效应的可能性是多少。所以这一课题牵涉到概率论的知识，但由于我们的数学基础还不够，我们不想过多地引用这方面的知识。

其次，从物理观点来说更重要的理由是，原子的实际行为并不遵循经典力学规律，而是遵循量子力学规律，因此，在我们了解量子力学以前，不可能得到对这个课题的正确的理解。这里，和弹子房弹球、汽车等情况不同，经典力学规律和量子力学规律之间的差别是非常重要和非常显著的，以致由经典物理学推导出来的许多结果从根本上来说就不正确。因此，这里学到的某些东西有一部分必须抛弃，然而，我们将指明每一种结果不正确的情况，以便了解它的“界限”何在。在前两章中讨论量子力学的理由之一就是给你们一个概念，使能了解为什么经典力学在许多方面或多或少是不正确的。

为什么我们现在就要处理这个课题呢？为什么不等上一年半载，直到我们更好地掌握了概率的数学理论，并且学了一点量子力学后，再以更为彻底的方式来处理它呢？回答是——这是一个困难的课题，学习它的最好方式是慢慢来！首先要做的是，使我们对不同场合下应当发生的情况多少获得一些概念，这样，以后当我们对这些规律了解得更清楚时，就能更好地用公式来表达它们。

任何一个企图分析实际问题中物质性质的人，可能都想从写出基本方程式出发，然后再从数学上求出它们的解。虽然有一些人试图采用这一条途径，但他们都是这个领域中的失败者；真正的成功来自那些从物理观点出发考虑的人，他们对于要做的事情具有大致的概念，并且在给定的复杂状况下知道哪些是重要的，哪些是次要的，然后开始作正确的近似。这些问题是如此复杂，即使获得那种不精确和不完全的初步理解也是很有价值的，因此在整个物理课程中，我们将一再碰到这个课题，而且一次比一次更为准确。

就在现在开始这个课题的另一个原因是，我们已经在例如化学上用过许多这方面概念，而且甚至在高中时，就听到过其中的某些内容。因此，了解这些事情的物理基础是颇有意义的。

举一个有趣的例子来说，我们都知道，在相同温度、相同压力下，相同体积的气体含有同等数目的分子。阿伏伽德罗把倍比定律，即当两种气体在化学反应中化合时，所需的体积总是成简单的整数比的定律，理解为相同体积中含有相同数目的原子。但是，为什么它们含有相同数目的原子？我们能从牛顿定律推出原子数目应该相等吗？本章中我们将谈到这种特殊的情况。在以下几章中，我们将讨论包括压强、体积、温度及热量的其他各种现象。

我们也发现，不从原子观点出发也同样能着手处理这个课题，并且在物质的性质上存在着许多相互联系。比如说，当我们压缩某个东西时，它会变热；如果把它加热，它就膨胀。在这两件事实之间存在着一种联系，它可以独立于更深一层的机制而推出。这个课题称为热力学。当然，对热力学的最深刻的理解来自对更深一层的实际机制的理解，而这就是我们将要做的：我们从一开始就采用原子的观点，并用它来了解物质的各种性质和热力学定律。

那么，下面就让我们从牛顿力学定律出发来讨论气体的性质。

\section{气体的压强}

首先，我们知道气体会产生压强。我们必须清楚地理解它是怎么产生的。如果我们的耳朵比现在灵敏几倍，就会听到持续的冲击噪声。进化论幸而没有使耳朵发展到这种地步，因为这样灵敏的耳朵毫无用处——不然我们将听到重复不停的吵闹声。原因是耳膜与空气相接触，而空气里有大量持续运动的分子，这些分子撞击在耳膜上。在撞击耳膜时，它们造成了无规则的咚、咚、咚的声音，这种声音我们听不见，因为原子非常小，以至于耳朵的灵敏度不足以感觉到。原子这样不停地撞击的结果是将耳膜推开，但是当然，在耳膜的另一边也有同样的原子在不停地撞击，因而作用在耳膜上的净力为零。如果我们从一边抽去空气，或者改变两边空气的相对数量，那么耳膜就将被推向这一边或那一边，因为在一边的撞击的量将大于另一边。当我们乘电梯或飞机时，由于上升得太快，特别是如果我们还患有重感冒的话，有时就会有这种不舒服的感觉（当我们感冒时，由于发炎而使通过咽喉联系耳膜内部空气与外部空气之间的导管关闭了，这样，内外压强就不能很快地保持平衡）。

为了定量地分析这种情况，我们设想有一个充满大量气体的容器，容器的一端是一个可移动的活塞（图\ref{figure:一个具有无摩擦活塞的容器中的气体原子}）。我们希望找出由于容器中有原子存在所引起的作用在活塞上的
\begin{wrapfigure}{r}{0.2\textwidth}
    \centering
    \includegraphics[width=0.19\textwidth]{Chapter39/一个具有无摩擦活塞的容器中的气体原子}
    \caption{一个具有无摩擦活塞的容器中的气体原子}
    \label{figure:一个具有无摩擦活塞的容器中的气体原子}
\end{wrapfigure}
力是多少。容器的体积为$V$，当原子在容器内以各种不同的速度来回运动时，它们就撞击在活塞上。假定在活塞外部没有什么东西，而是真空，那么它会怎么样呢？如果活塞是自由的，没有什么东西顶着，那么它每次受到撞击时，都将得到一点点动量，于是便会渐渐地从容器中被推出。所以为了使它不从容器中被推出，我们必须加上一个力$F$。问题在于要用多大的力？表示力的一种方法是考察每单位面积上的力：如果$A$是活塞的面积，那么作用在活塞上的力可以写成某个数乘以面积。于是，我们定义压强为加在活塞上的力除以活塞的面积
\begin{equation}
    \label{Eq:I:39:1}
    P = \frac{F}{A}
\end{equation}
为了更好地理解这个概念（为其他目的我们也曾推导出它），必须注意由于使活塞移动距离（$-dx$）压缩气体时，对气体所作的元功$dW$，应当是力乘以压缩的距离，而按照式（\ref{Eq:I:39:1}），就是压强乘以面积再乘以距离，即压强乘体积变化的负值
\begin{equation}
    \label{Eq:I:39:2}
    dW = F(-dx) = -PAdx = -PdV
\end{equation}
（面积$A$乘以距离$dx$等于体积的变化）。这里的负号是因为气体受压缩时，它的体积减小；如果考虑到这一点，就可以看出在气体受到压缩时，外界对它做了功。

我们必须施加多大的力来平衡分子的碰撞呢？活塞从每次碰撞中获得一定的动量。每秒钟都有一定的动量倾注在活塞上，使它开始运动。为了使活塞保持不动，必须从我们的力中每秒钟输回给它同样大小的动量。当然，这个力正是我们每秒钟必须输送给它的动量的数量。还可以用另一种方法来考虑这个问题：如果让活塞自由活动，它将由于碰撞而获得速率；而每一次碰撞都使速率增大一点，这样它就不断地受到加速。活塞的速率的增加，或者说加速度正比于作用在它上面的力。于是我们看到这个力等于由于分子\footnote{原文如此，在本节文字中，作者有时似将分子和原子这两个概念混用。——译者注}的碰撞每秒钟传给活塞的动量，而我们前面已经谈到过它等于压强乘以面积。

计算每秒钟施加给活塞的动量是不难的，我们可以分两步来求：首先，找出一个特定的原子在和活塞的一次碰撞中传给活塞的动量，然后，再乘上每秒钟原子与活塞壁发生碰撞的次数。力就是这两个因子的乘积。现在我们来看看两个因子是怎样的：首先，我们假定活塞对原子来说是一个理想的“反射体”，否则，整个理论就是错误的，活塞将开始变热，事情将发生变化。但是最后，当平衡建立后，总的效果是碰撞仍是有效的完全弹性碰撞。平均而言，每个粒子飞来和离开时都带有相同的能量。我们想象气体处在稳定的状态下，并且由于活塞静止不动，所以对于活塞我们没有损耗能量。在这种情况下，如果某一质量的粒子以一定的速率飞来，它也以同样的质量和同样的速率离开。

如果原子的速度为$\boldsymbol{v}$，$\boldsymbol{v}$的$x$分量为$v_x$，则“入射”粒子动量的$x$分量为$mv_x$；“出射”粒子动量的分量和它相等，这样，粒子在一次碰撞中施加给活塞的总动量是$2mv_x$，因为这个粒子是被“反射”回来的。

现在，我们要知道每秒钟内原子与活塞的碰撞次数；或者说$dt$时间内的碰撞次数，然后再除以$dt$。有多少原子打上去呢？我们假设在体积$V$中有$N$个原子，即单位体积内的原子数为$n=N/V$。为了找出有多少个原子打在活塞上，我们注意到，给定一个时间$t$，若某个粒子离活塞足够近，并且具有一定的指向活塞的速度，就能在这段时间$t$内碰上活塞。如果它离活塞太远，那么在时间$t$内只能向着活塞跑过一段路程，而不能到达活塞。非常清楚，只有离活塞的距离在$v_xt$之内的那些分子，才会在时间$t$内打在活塞上。因而在时间$t$内的碰撞次数等于在距离为$v_xt$内的区域里的原子数，并且由于活塞的面积是$A$，所以能在时间$t$内碰到活塞的原子所占有的体积是$v_xtA$。但是碰撞到活塞的原子数等于这个体积乘以单位体积的原子数，即为$nv_xtA$。当然，我们并不要求时间$t$内的碰撞次数，而是每秒钟的碰撞次数，因而只要除以时间$t$，就得到$nv_xA$（这个时间$t$可以很短；如果我们觉得需要更精确一些，可以称它为$dt$，然后再求微商，但结果是相同的）。

这样，我们求得力为
\begin{equation}
    \label{Eq:I:39:3}
    F=nv_xA\cdot 2mv_x
\end{equation}
可见，若保持粒子密度不变，而改变面积，则力与面积成正比。于是压强为
\begin{equation}
    \label{Eq:I:39:4}
    P = 2nmv_x^2
\end{equation}

不过我们注意到在这个分析中有一些小小的麻烦：首先，并非所有分子都具有同样的速度，而且它们并不都以同样的方向运动。因而，所有这些$v_x^2$项都不相同！因为每个分子都有它自己的贡献，所以，我们要做的自然是对这些$v_x^2$取平均。即我们要求的是$v_x$平方对所有分子的平均值
\begin{equation}
    \label{Eq:I:39:5}
    P = nm\langle v_x^2 \rangle
\end{equation}
这里我们是不是忘了写上因子$2$？不，在所有的原子中，只有一半是朝着活塞跑的，另一半朝着相反的方向运动。对正的$v_x$，$v_x^2$的平均值等于对所有$v_x$所求平均值的一半。

当原子向四面八方运动时，显然在“$x$方向”上没有什么特殊之处；原子同样可以上下、左右、前后地运动。因此在运动过程中，表征原子在一个方向上平均运动的$\langle v_x^2 \rangle$值和在其他两个方向上的平均值全都相等
\begin{equation}
    \label{Eq:I:39:6}
    \langle v_x^2 \rangle = \langle v_y^2 \rangle = \langle v_z^2 \rangle 
\end{equation}
因此，只要稍微用一点点数学技巧就可看出，每一项等于这三项总和的三分之一，而这三项之和当然就是速率的平方
\begin{equation}
    \label{Eq:I:39:7}
    \langle v_x^2 \rangle = \frac{1}{3}\langle v_x^2 + v_y^2 + v_z^2 \rangle = \frac{\langle v^2 \rangle}{3}
\end{equation}
这个式子的方便之处在于我们毋须考虑任何特殊的方向，于是，可把压强公式重新写为
\begin{equation}
    \label{Eq:I:39:8}
    P = \left ( \frac{2}{3} \right )n \left \langle \frac{mv^2}{2} \right \rangle
\end{equation}
把最后一个因子写成$\langle mv^2/2 \rangle$是因为这是一个分子的质心运动的动能。由此我们得出
\begin{equation}
    \label{Eq:I:39:9}
    PV = N \left ( \frac{2}{3} \right )\left \langle \frac{mv^2}{2} \right \rangle
\end{equation}
如果我们知道速率，用这个公式就可计算出压强有多大。

举一个十分简单的例子，我们考虑氦气或其他任何气体，例如水银蒸气，足够高温度下的钾蒸气或氩气，其中所有分子都是单原子分子，对这些单原子分子，我们可以假定在分子中没有内部运动。如果是复杂的分子，其中就可能存在某些内部运动，相互间的振动，等等。假设我们忽略这些运动（实际上，这些运动是很重要的，以后我们再回过头来考虑）；这样做还是可行的。由于我们假定原子内部的运动可以不考虑，质心运动的动能就是分子所具有的全部能量。所以，对于单原子气体，动能就是总能量。通常我们称$U$为总能量（有时也称$U$为总内能。我们或许会感到奇怪，因为对气体来说，并没有外能），它就是在气体中，或无论什么东西中所有分子的全部能量。

我们假设单原子气体的总能量$U$等于原子数乘以每个原子的平均动能，因为我们不考虑原子内部任何可能的运动或激发。于是，在这些条件下，我们有
\begin{equation}
    \label{Eq:I:39:10}
    PV = \frac{2}{3}U
\end{equation}

附带提一下，我们可以在这里停一下并找出下述问题的答案：假如我们缓慢地压缩盛在一个容器中的气体，要把体积压缩需要多大的压强？这不难求出，因为压强是总能量的三分之二再除以$V$。当我们压缩气体时，就对气体做功，因而增加了能量$U$。于是，我们就得到了某种微分方程：如果我们从具有一定内能和一定体积的已知条件出发，那么就能求出压强。一旦我们开始压缩气体，内能就增加，同时体积减小，从而压强增大。

这样，我们就需要求解微分方程。我们马上就来解它。不过首先必须强调指出，在压缩这一气体时，我们假设所有的功都转变为增加气体内部原子的能量。我们可能会问：“这种假设是否有必要？它所做的功难道还会跑到别的什么地方去吗？”结果表明，它是能够跑到别的地方去的。有一种所谓通过器壁“漏热”的现象：当热的（即快速运动的）原子撞击器壁时，加热器壁，使能量跑出去。现在我们假定不发生这种情况。

为了略微更普遍一些，虽则我们仍然对气体作某些十分特殊的假设，我们将不写$PV=2U/3$，而写
\begin{equation}
    \label{Eq:I:39:11}
    PV = (\gamma - 1)U
\end{equation}
按照习惯，把它写成$(\gamma-1)$乘以$U$，因为在以后处理的少数其他情况中，U前面的因子不是$2/3$，而是其他的数值。所以，为了更一般起见，我们称它为$(\gamma-1)$，因为人们已经这样称呼了近一百年了。对于像氦这样的单原子气体，因为$(5/3-1)$是$2/3$，所以$\gamma$等于$5/3$。

我们已经注意到，压缩气体时所做的功是$-PdV$。既不加入热能也不取走热能的压缩过程称为绝热（adiabatic）压缩，它是由希腊字a（不）+ dia（穿）+ bainein（过）而来的。（“绝热”这个词在物理中有几种不同的用法，有时很难看出它们之间有什么共同含义。）这就是说，对于绝热压缩，所作的全部功都转变为内能。没有其他能量损失——这就是关键，因而我们有$PdV=-dU$。但因$U=PV/(\gamma-1)$，可得
\begin{equation}
    \label{Eq:I:39:12}
    dU = \frac{PdV + VdP}{\gamma - 1}
\end{equation}
因而我们有$PdV=-(PdV+VdP)/(\gamma-1)$，或者整理一下，得
\begin{equation*}
    \gamma PdV = -VdP
\end{equation*}
或
\begin{equation}
    \label{Eq:I:39:13}
    \frac{\gamma dV}{V} + \frac{dP}{P} = 0
\end{equation}
幸运的是，假定$\gamma$是常数（比如对单原子气体），我们就能进行积分：即$\gamma\ln V+\ln P=\ln C$，这里$C$是积分常数。对两边取指数，就得到定律
\begin{equation}
    \label{Eq:I:39:14}
    PV^\gamma = C~~(\text{常数})
\end{equation}
换句话说，在绝热条件下，在气体的压缩过程中，由于没有热量流失，因而温度升高，对单原子气体来说，压强乘以体积的$5/3$次方是一个常数！虽然我们是从理论上推出这个结论的，但事实上，单原子气体在实验上的表现也是如此。

\section{辐射的压缩性}

我们可以举另一个气体分子动理论的例子，它在化学上用得不太多，但在天文学上有用。设有一个装有大量光子的容器，其中温度非常高（当然，高热恒星中的气体就可看成是这样的容器。太阳还不够热；在那里仍有太多的原子，但是在某种更热的恒星中，会有更高的温度，我们可以忽略原子，而假设在容器中唯一的客体是光子）。而光子具有一定的动量$\boldsymbol{p}$（当我们学习分子动理论时，会一再遇到麻烦：$p$既是压强，又是动量；$v$既是体积，又是速度；$T$既是温度，又是动能、时间或者力矩；我们必须保持警惕）。这里$\boldsymbol{p}$是动量，是矢量。按照与前面相同的分析，正是矢量$\boldsymbol{p}$的$x$分量产生“反冲”的，矢量$\boldsymbol{p}$的$x$分量的两倍是在反冲中给出的动量。于是$2\boldsymbol{p}_x$代替了$2mv_x$，而在计算碰撞次数时，$v_x$仍为$v_x$，这样当我们继续采取以前的所有步骤后，发现式（\ref{Eq:I:39:4}）中的压强可用下式来代替
\begin{equation}
    \label{Eq:I:39:15}
    P = 2np_xv_x
\end{equation}
在作平均时，它变为$n$乘以$p_xv_x$的平均值（因子$2$的情况同上），最后计入另外两个方向，我们求得
\begin{equation}
    \label{Eq:I:39:16}
    PV = \frac{N}{3}\langle \boldsymbol{p} \cdot \boldsymbol{v} \rangle
\end{equation}
此式与（\ref{Eq:I:39:9}）相符，因为动量是$m\boldsymbol{v}$；只是它稍微更一般些，如此而已。总之，压强乘体积等于原子总数乘$(\boldsymbol{p}\cdot\boldsymbol{v})/3$的平均值。

对光子来说，$\boldsymbol{p}\cdot\boldsymbol{v}$是什么？动量与速度方向相同，而速度就是光速，所以这就是每个光子的动量与光速的乘积。每个光子的动量乘光速是它的能量：$E=pc$，因而这些项就是每个光子的能量，当然，我们应当取光子的平均能量乘光子数。这样，$PV$的乘积是气体中能量的三分之一
\begin{equation}
    \label{Eq:I:39:17}
    PV = \frac{U}{3}(\text{光子气体})
\end{equation}
于是，对光子来说，由于前面的系数是$1/3$，即在式（\ref{Eq:I:39:11}）中的$(\gamma-1)$是$1/3$，所以$\gamma=4/3$，因而我们发现容器内的辐射满足规律
\begin{equation}
    \label{Eq:I:39:18}
    PV^{4/3} = C~~(\text{常数})
\end{equation}
从而我们知道了辐射的压缩性！这就是用在分析恒星上辐射压强贡献的关系式。这也是我们算出辐射压强的方法，它也表示当我们压缩光子气体时，压强是怎样变化的。瞧，多么奇妙的事情我们也都能处理！

\section{温度和动能}

到现在为止我们还没有考虑温度；这个概念是我们有意回避的。在压缩气体时，分子的能量增加，我们通常就说气体变热了；现在我们想了解它和温度的关系。如果要做一个所谓等温而不是绝热的实验，该怎么办呢？我们知道，如果让盛有气体的两个容器彼此紧贴足够长时间后，甚至即使开始时它们处在我们称之为不同的温度之下，最后也将具有相同的温度。这意味着什么呢？这意味着它们将达到的状态就是如果我们让它们单独存在足够长时间后它们所应达到的那种状态！我们所说的相同温度指的正是当物体放在一起相互作用足够长时间后达到的最终状态。

如图\ref{figure:两种不同的单原子气体的原子被一可移动的活塞隔开}所示，假设在一个被可移动活塞分开的容器里放有两种气体，我们现在来
\begin{wrapfigure}{r}{0.25\textwidth}
    \centering
    \includegraphics[width=0.24\textwidth]{Chapter39/两种不同的单原子气体的原子被一可移动的活塞隔开}
    \caption{两种不同的单原子气体的原子被一可移动的活塞隔开}
    \label{figure:两种不同的单原子气体的原子被一可移动的活塞隔开}
\end{wrapfigure}
考虑将会出现什么情况（为了简单起见，假定这两种气体都是单原子气体，如氦与氖）。在容器（1）中的原子质量为$m_1$，速度为$v_1$，单位体积原子数为$n_1$，而在另一个容器（2）中，原子质量为$m_2$，速度为$v_2$，单位体积原子数为$n_2$。它们的平衡条件是什么？

很明显，从左边来的碰撞必定会使活塞向右移动，并且压缩另一边的气体，使它的压强升高，从而活塞来回移动，逐渐停在某个使两边压强相等的位置上。所以我们可以使压强相等；而这正意味着单位体积的内能相等，或者说每一边的单位体积原子数$n$与平均动能的乘积彼此相等。最后，我们想要证明的是，这些数目本身也相等。到目前为止我们所知道的只是，根据式（\ref{Eq:I:39:8}），由于压强相等，所以原子数乘平均动能相等
\begin{equation*}
    n_1 \left \langle \frac{m_1v_1^2}{2} \right \rangle = n_2 \left \langle \frac{m_2v_2^2}{2} \right \rangle
\end{equation*}
我们必须了解到这并不是经过长时间后所达到的唯一的条件，当相应于温度相等的真正完全的平衡建立起来时，还必将发生某些更缓慢的过程。

为了搞清楚这种概念，假设作用在左边的压强是在高密度和低速度的情况下产生的。具有大的$n$和小的$v$，与具有小的$n$与大的$v$可以得出同样的压强。原子可能运动得很慢，但挤得很紧密，也可能密度较小，但撞击力很强。它能永久地像这样继续下去吗？一开始我们可能这样想，但是再思考一下就会发现我们忘掉了一件重要的事情。也就是说，中间的活塞不再受到一种稳定的压力；正像我们一开始就讲过的人的耳膜那样，因为撞击不是绝对均匀，它左右摆动。结果没有获得持续、稳定的压力，而是不断地敲击——压强在变化，因而活塞轻轻地晃动。假定右边的原子晃动得不太大，而左边的原子数较少，原子间隔较远，但带有较大的能量，那么，活塞将会从左边得到一个较大的冲量，因而将驱动右边缓慢运动的原子，使它们获得更大的速率（每当原子与活塞碰撞时，它不是得到就是失去能量，取决于原子敲打在活塞上时活塞向哪个方向运动）。这样，由于碰撞的结果，活塞本身就反复地晃动，而这会使其他气体震动——它将能量传给了另一些原子，使它们运动得更快，直到与活塞所引起的晃动相平衡为止。当活塞以这样的方均速率运动，使它在单位时间内从原子获得的能量和它送回给原子的能量的比率大致相等时，系统就达到某种平衡。这样，活塞在速率上就获得一定的平均不规则性，这正是我们要找的。当我们找到它以后，就能更好地解决问题，因为这些气体将调节它们的速度，直到在单位时间内通过活塞彼此交换的能量变成相等为止。

在这种特殊情况下，要描述活塞运动的细节是十分困难的；虽然理解起来非常简单，但分析起来却比较困难。在我们分析这个问题以前，我们先来分析另一个问题：设有一个容器，其中包含由两种不同分子组成的气体，两种分子的质量分别为$m_1$及$m_2$，速度为$v_1$及$v_2$，等等；现在就有了一种更为密切的关系。如果所有的第二种分子都静止不动，那么这种情况将不会延续下去，因为它们受到第一种分子的碰撞，从而获得速度。如果所有的第二种分子都运动得比第一种分子快，那么这种情况也不会持续多久，它们将反过来把能量传给第一种分子。所以当在同一容器中存在两类气体时，问题就是要求出确定两者之间的相对速度的规则。

这仍然是一个十分困难的问题，我们将这样来解决。首先，我们考虑下面一个附属问题（这又是那种情况之一——别管推导——最终的结果很容易记住，但推导方式却是十分巧妙的）。假设有两个质量不同的分子发生碰撞，我们在质心（CM）系来考察这个
\begin{wrapfigure}{l}{0.25\textwidth}
    \centering
    \includegraphics[width=0.24\textwidth]{Chapter39/从质心系来看两个不同分子间的碰撞}
    \caption{从质心系来看两个不同分子间的碰撞}
    \label{figure:从质心系来看两个不同分子间的碰撞}
\end{wrapfigure}
碰撞。为了避免复杂性，我们在质心观察碰撞。由碰撞过程中动量和能量守恒定律可知，两个分子碰撞后，它们运动的唯一可能方式是各自保持原来的速率，即只改变它们的方向。所以，我们就有了一种看来像图\ref{figure:从质心系来看两个不同分子间的碰撞}所示的一般碰撞。假设我们暂时从静止的质心系观察所有这类碰撞，并且设想这些分子起初全都沿水平方向运动。当然，在第一次碰撞后，其中有一些以某一角度运动。换句话说，如果起初它们全都沿水平方向运动，那么以后至少会有一些沿垂直方向运动。而在某些其他碰撞中，它们可能来自其他方向，然后又以另外的角度运动。因而即使在开始时分子都排成一个方向运动，它们也会向所有的角度飞离，飞离的分子还将飞离更多次，如此继续下去。最终将怎样分布呢？答案是：任何一对分子沿空间任何方向运动的可能性是相等的。此后进行的碰撞也不会改变这个分布。

分子沿所有方向运动的机会是相等的，不过我们怎样来表示这一点呢？当然，它们不可能沿某一特定方向运动，因为某一特定方向过于严格，所以我们必须说每单位“某某”。我们的概念是在以碰撞点为中心的球面上，通过任何一块面积的分子数正好等于通过在球面上任何其他相等面积的分子数。因而碰撞的结果将使分子的方向这样分布，使得球面的每个相等的面积有相同的分子通过。

附带说一下，如果我们只要讨论原来的方向和某个与它成角度$\theta$的其他方向，那么，有趣的是单位半径球的面积元为$\sin\theta d\theta$乘以$2\pi$（图\ref{figure:球截形的侧面积为}），这正是$-\cos\theta$的微分乘以$2\pi$。它意味着任意两个方向间夹角$\theta$的余弦取$-1$至$+1$间的任意值的可能性是相同的。

下面我们来考虑实际情况，这时不是质心系中的碰撞，而是两个速度矢量分别为$\boldsymbol{v}_1$与$\boldsymbol{v}_2$的原子跑到一起的碰撞。这时会发生什么情况呢？我们可以对这个带有速度矢量$\boldsymbol{v}_1$和$\boldsymbol{v}_2$的碰撞分析如下：

首先，有一个质心，质心的速度是权重和质量成正比的“平均”速度，即质心的速度是$\boldsymbol{v}_{CM}=(m_1\boldsymbol{v}_1+m_2\boldsymbol{v}_2)/(m_1+m_2)$。如果我们在质心系中观察这个碰撞，那么我们看到的碰撞正如图\ref{figure:从质心系来看两个不同分子间的碰撞}所示，彼此以一定的相对速度$\boldsymbol{w}$靠近。相对速度正好是$\boldsymbol{v}_1-\boldsymbol{v}_2$。现在，我们的想法是，首先，整个质心在运动，在质心系中有一个相对速度$\boldsymbol{w}$，接着，分子碰撞后按某个新方向离开。所有这些都是在质心系保持原来运动时发生的，没有任何变化。

那么，由此而产生的分布是怎样的呢？由前面的论证我们断定：在平衡时，相对于质心的运动方向来说，$\boldsymbol{w}$在一切方向的可能性相同\footnote{这个曾为麦克斯韦所采用的论证含有某些微妙之处。虽然结论是正确的，但这个结果并不能纯粹从我们前面用过的对称性的考虑得出，因为过渡到在气体中运动的参考系时，我们可以得到一个改变了的速度分布。但我们还未找到这个结果的简单证明。}。结果，相对速度的方向与质心运动方向之间没有任何特殊联系。当然，如果有，碰撞也会破坏这种联系，直到全部破坏掉为止，因而$\boldsymbol{w}$与$\boldsymbol{v}_{CM}$之间夹角之余弦的平均值是零，即
\begin{equation}
    \label{Eq:I:39:19}
    \langle \boldsymbol{w}\cdot\boldsymbol{v}_{CM} \rangle = 0
\end{equation}
但$\boldsymbol{w}\cdot\boldsymbol{v}_{CM}$可用$\boldsymbol{v}_1$及$\boldsymbol{v}_2$表示为
\begin{equation}
    \label{Eq:I:39:20}
    \boldsymbol{w}\cdot\boldsymbol{v}_{CM} = \frac{(\boldsymbol{v}_1 - \boldsymbol{v}_2)\cdot(m_1\boldsymbol{v}_1 + m_2\boldsymbol{v}_2)}{m_1 + m_2} = \frac{(m_1v_1^2 - m_2v_2^2) + (m_2 - m_1)(\boldsymbol{v}_1 \cdot \boldsymbol{v}_2)}{m_1 + m_2}
\end{equation}

我们看看$\boldsymbol{v}_1 \cdot\boldsymbol{v}_2$这项；$\boldsymbol{v}_1 \cdot\boldsymbol{v}_2$的平均值是什么？也就是说，一个分子的速度在另一个分子的速度方向上的分量的平均值是多少？显然，找到以一种方式运动的任何给定的分子的可能性与找到以另一种方式运动的给定的分子的可能性相同。速度$\boldsymbol{v}_2$沿任意方向的平均值是零。所以，在$\boldsymbol{v}_1$方向上，$\boldsymbol{v}_2$的平均值显然为零。因而$\boldsymbol{v}_1 \cdot\boldsymbol{v}_2$的平均值就是零！由此能推知$m_1v_1^2$的平均值必须等于$m_2v_2^2$的平均值。这就是说，两个分子的平均动能必定相等
\begin{equation}
    \label{Eq:I:39:21}
    \frac{1}{2}m_1v_1^2 = \frac{1}{2}m_2v_2^2
\end{equation}
如果在气体中有两种原子，那么可以证明，而且我们相信已经证明当它们两者都处在同样气体、同样容器内，并且处于平衡状态时，则一种原子的平均动能与另一种原子的平均动能相等。这意味着，重的原子将比轻的原子运动得慢一些；这不难用气垫中不同质量的“原子”的实验证明。

\begin{wrapfigure}{r}{0.25\textwidth}
    \centering
    \includegraphics[width=0.24\textwidth]{Chapter39/两种气体在一个具有半透膜的容器中}
    \caption{\footnotesize 两种气体在一个具有半透膜的容器中}
    \label{figure:两种气体在一个具有半透膜的容器中}
\end{wrapfigure}
现在我们再进一步断言，如果在一个容器内有两种不同的分开的气体，当它们最后达到平衡时，即使不在同一容器内，它们也将具有相同的平均动能。我们可以用几种方法来论证。一种论证方法假定在容器内有一块固定的上面开有一个小孔的隔板（图\ref{figure:两种气体在一个具有半透膜的容器中}），使得一种气体可以通过小孔漏出去，而另一种气体则因为分子太大而不能漏出。当达到平衡时，我们知道，在混合的那部分气体中，它们具有同样的平均动能，而通过小孔的那些分子没有失去动能，因而在纯气体中的平均动能和在混合气体中的平均动能必然相等。这个论证还不太令人满意，因为可能不存在把一种气体分子与另一种气体分子分开的小孔。

现在我们回到活塞的问题上来。我们可以提出一个论证来说明这个活塞的动能必定也是$m_2v_2^2/2$，实际上，这就是由活塞作纯粹水平运动引起的动能，这样，如果忽略它的上下运动，它所具有的动能同样必然是$m_2v_{2x}^2/2$。类似地，从另一边的平衡条件，又可证明活塞的动能是$m_2v_{1x}^2/2$。即使当活塞并不处于气体中央，而是在气体的一边时，虽然证明略为困难一些，我们仍然可以作出同样的论证，也就是由于所有碰撞的结果，活塞的平均动能与气体分子的平均动能彼此相等。

如果这样做还不满意，可以设想一种人为的例子，认为平衡是由一个能打到所有各边的物体产生的。假定我们有一根每一端都有一个小球的短棒，它穿过活塞插在一个无摩擦的滑动万向接头上。每个小球都像一个分子那样是圆的，可以接受所有方向上的撞击。整个物体的总质量是$m$。现在，和之前一样，气体分子的质量各为$m_1$与$m_2$。用前面的分析可知，碰撞的结果是，由于受一边分子的撞击，$m$的动能平均说来是$m_1v_1^2/2$。类似地，由于受另一边分子的撞击，$m$的动能平均说来必然是$m_2v_2^2/2$。所以，当它们处于热平衡时，两边必定有相同的动能。因而，虽然我们只对混合气体证明了这一点，但很容易推广到在同样温度下的两种不同的、分离的气体中去。

这样，当两种气体处于相同温度时，质心运动的平均动能相等。

分子的平均动能只是“温度”的一种特性，由于它只是“温度”的特性，而不是气体的特性，因而我们可以利用它作为温度的定义。于是分子的平均动能即为温度的某种函数。但是，谁来告诉我们温度该用什么尺度呢？我们自己可以任意地定义温度的尺度，使得平均动能和温度成线性正比关系。要做到这一点最好的办法就是把平均动能本身叫做“温度”。这是一个最简单的函数。遗憾的是，温度的尺度已经按其他方式选定了，所以我们不直接把平均动能称为温度，而是在分子的平均动能与所谓开尔文绝对温度一度之间加上一个常数的转换因子。比例常数是$k=1.38\times 10^{-23}~J/K$\footnote{摄氏温标与开尔文温标一样，但$0~^\circ\text{C}$选在$273.16~K$，所以$T=273.16+$摄氏温度。}。因此，如果$T$是绝对温度，按我们的定义，分子的平均动能是$3kT/2$（$3/2$是为了方便而引入的，以便在其他地方去掉它）。

我们指出和运动沿任何特定方向的分量相联系的动能是$kT/2$。因为平均动能包含三个独立的运动方向，所以总和为$3kT/2$。

\section{理想气体定律}

当然，现在我们可以把温度的定义代入式（\ref{Eq:I:39:9}）中，从而找到气体压强与温度之间的函数关系：即压强乘体积等于原子总数乘以普适常数$k$再乘以温度
\begin{equation}
    \label{Eq:I:39:22}
    PV = NkT
\end{equation}
而且，在同样的温度、同样的压强与同样的体积下，原子数是确定的；这也是一个普适常数！所以，根据牛顿定律，在同样的温度和同样的压强下，相同体积的不同气体中具有相同的分子数。这是一个令人惊异的结论！

实际上在处理分子时，因为分子数量太大，化学家人为地选择了一个特定的很大的数目，而给它一个名称，称为1摩尔（mol），$1~mol$只是个方便的数。为什么他们不选$10^24$个分子这样的偶数，这只是个历史问题而已。化学家为了方便起见恰巧选择了标准情况下的分子数$N_0=6.02\times 10^{23}$，而称它为$1~mol$的分子数。所以化学家并不是以分子为单位来测量分子数，而是以摩尔数\footnote{化学家所说的分子量是指$1~mol$的分子以克计量的质量。$1~mol$是这样定义的，即要使1 mol碳同位素12（也就是说，原子核内有六个质子和六个中子）的质量正好是12克。}来测量的。根据$N_0$可以写出摩尔数，乘上$1~mol$中原子数，再乘以$kT$，而且如果我们需要的话，可以取$1~mol$的原子数乘以$k$，那就是$1~mol$的$k$值，而把它叫做其他的什么，我们叫它$R$。$1~mol$的$k$是$8.317~J:~R=N_0k=8.317~J\cdot mol^{-1}\cdot K^{-1}$。这样我们也发现气体定律可写为摩尔数（也称为$N$）乘以$RT$，或原子数乘$kT$
\begin{equation}
    \label{Eq:I:39:23}
    PV = NRT
\end{equation}
它们是完全一样的，只是测量数目的尺度不同而已。我们用1作为单位，而化学家用$6\times 10^{23}$作为单位！

现在我们对气体定律再作一点说明，这与非单原子分子组成的气体的定律有关。我们只处理了单原子气体的原子的质心运动。如果还有一些力存在，会出现什么情况呢？首先，考虑活塞被一水平方向的弹簧拉住的情况，这时，有力作用于活塞上。当然，在任何时刻，原子与活塞间的无规则晃动的交换都和这时的活塞位置无关。平衡条件是相同的。无论活塞在哪里，它的运动速率都能刚好正确地将能量传递给分子。弹簧的存在与否并没有任何不同。平均说来，活塞运动应取的速率是相同的。因而，在一个方向上动能平均值是$kT/2$的这个定理，无论有无力存在，都是正确的。

例如，考虑由原子$m_A$和$m_B$组成的双原子分子的情况。我们已经证明了$A$部分的质心运动与$B$部分的质心运动有下列关系
\begin{equation*}
    \left \langle \frac{1}{2}m_Av_A^2 \right \rangle = \left \langle \frac{1}{2}m_Bv_B^2 \right \rangle = \frac{3}{2}kT
\end{equation*}
如果它们合在一起后，是否也会如此呢？虽然它们合在一起，当它们在那里不断自转和旋转，当别的分子撞击它们，和它们交换能量时，唯一要计入的因素是分子跑得多快。只有这一点才确定了它们在碰撞中交换的能量有多快。在那一瞬间，力不是主要的，因此即使有力存在，同样的原理也是正确的。

最后，我们证明，不考虑分子内部运动时，气体定律同样成立。实际上以前我们并没有包括内部运动，只研究了单原子气体。但现在我们将证明，如果把整个系统考虑成一个总质量为$M$，具有质心速度的单个物体的话，则
\begin{equation}
    \label{Eq:I:39:24}
    \frac{1}{2}Mv_{CM}^2 = \frac{3}{2}kT
\end{equation}
换句话说，我们既可以考虑分开的部分，也可以考虑整体！原因在于：双原子分子的质量是$M=m_A+m_B$，而质心速度等于$\boldsymbol{v}_{CM}=(m_A\boldsymbol{v}_A+m_B\boldsymbol{v}_B)/M$。现在要求出$\langle v^2_{CM} \rangle$。取$\boldsymbol{v}_{CM}$的平方，得
\begin{equation*}
    v_{CM}^2 = \frac{m_A^2v_A^2 + 2m_Am_B\boldsymbol{v}_A\cdot\boldsymbol{v}_B + m_B^2v_B^2}{M^2}
\end{equation*}
在两边各乘以$M/2$，然后取平均，可得
\begin{equation*}
    \begin{aligned}
        \frac{1}{2}Mv_{CM}^2 &= \frac{m_A\dfrac{3}{2}kT + 2m_Am_B\langle \boldsymbol{v}_A\cdot\boldsymbol{v}_B \rangle + m_B\dfrac{3}{2}kT}{M} \\
        &= \frac{3}{2}kT + \frac{2m_Am_B\langle \boldsymbol{v}_A\cdot\boldsymbol{v}_B \rangle}{M}
    \end{aligned}
\end{equation*}
[我们用到了$(m_A+m_B)/M=1$]。现在，$\langle \boldsymbol{v}_A\cdot\boldsymbol{v}_B \rangle$是多少（最好是零）？为了求出结果，我们利用原先的假定，即相对速度$\boldsymbol{w}=\boldsymbol{v}_A-\boldsymbol{v}_B$不会特别偏向于哪个方向，也就是说，它在任何方向上的平均分量均为零，即
\begin{equation*}
    \langle \boldsymbol{w}\cdot\boldsymbol{v}_{CM} \rangle = 0
\end{equation*}
但$\boldsymbol{w}\cdot\boldsymbol{v}_{CM}$是什么？它是
\begin{equation*}
    \begin{aligned}
        \boldsymbol{w}\cdot\boldsymbol{v}_{CM} &= \frac{(\boldsymbol{v}_A - \boldsymbol{v}_B) \cdot (m_A\boldsymbol{v}_A + m_B\boldsymbol{v}_B)}{M} \\
        &= \frac{m_Av_A^2 + (m_B-m_A)(\boldsymbol{v}_A\cdot\boldsymbol{v}_B) - m_Bv_B^2}{M}
    \end{aligned}
\end{equation*}
因此，由于$\langle m_Av_A^2 \rangle = \langle m_Bv_B^2 \rangle$，平均后将第一项与最后一项消去，只剩下
\begin{equation*}
    (m_A - m_B) \langle \boldsymbol{v}_A\cdot\boldsymbol{v}_B \rangle = 0
\end{equation*}
如果$m_A \neq m_B$，就可得出$ \langle \boldsymbol{v}_A\cdot\boldsymbol{v}_B \rangle = 0$，因此，将整个分子的总体运动考虑成一个质量为$M$的单个粒子运动时，它所具有的平均动能等于$3kT/2$。

我们附带也证明了，如果不考虑质心的整体运动，则双原子分子内部运动的平均动能也是$3kT/2$！因为分子各部分的总动能是$m_Av_A^2/2 + m_Bv_B^2/2$，而其平均值是$3kT/2+3kT/2=3kT$。而质心运动的平均动能是$3kT/2$，因而分子中两个原子的转动与振动的平均动能是它们之差，也等于$3kT/2$。

关于质心运动平均动能的定理是普遍的：当把任何物体考虑为一个整体时，无论是否有力存在，在这个物体的每个独立的运动方向上，其平均动能都是$kT/2$。这些“独立的运动方向”有时也称为系统的自由度。由$\gamma$个原子组成的分子的自由度数是$3\gamma$，因为每个原子都需要有三个坐标来确定它的位置。分子的总动能既可以表示为各个原子的动能的和，也可以表示为质心运动的动能与内部运动的动能之和，后者有时可以表示为分子转动动能与振动动能之和，但这是一个近似。把我们的定理应用到$\gamma$个原子的分子时表明，分子平均动能将是$3\gamma kT/2$，其中$3kT/2$是整个分子的质心运动的动能，其余的$3(\gamma-1)kT/2$则是分子内部的振动与转动动能。
